\subsection{Loans}
\content{
        Formulas so far: \\
        \fbox{\begin{minipage}{\textwidth}
        For these formulas, we have the following: 
        $P_0$ is the principal. 
        $r$ is the interest rate. 
        $t$ is the number of time intervals (specified by problem) which
        have passed. 
        $A$ is the amount paid back.
        $I$ is the interest.

        \begin{center}
        \begin{tabular}{|c|c|}
        \hline
        \T\B Simple Interest &  $ A = P_0(1+r) $ \\
        \hline
        \T\B Simple Interest over Time & $  A = P_0(1+rt) $ \\
        \hline
        \end{tabular}
        \end{center}
        \end{minipage}}
        \fbox{\begin{minipage}{\textwidth}
        For these formulas, we have the following: 
        $P_0$ is the principal. 
        $r$ is the annual interest rate.
        $t$ is the number of years which have passed. 
        $A$ is the total amount.
        $k$ is the number of times interest is compounded per year. 
        $d$ is the amount we deposit (or withdraw) from the account regularly.

        \begin{center}
        \begin{tabular}{|c|c|}
        \hline
        \T\B Compound Interest & $ A =
        P_0\left(1+\frac{r}{k}\right)^{kt} $ \\
        \hline
        \T\B Savings Annuity & $ A =
        d\frac{\left(\left(1+\frac{r}{k}\right)^{kt}-1\right)}{\frac{r}{k}} $\\
        \hline
        \T\B Payout Annuity & $P_0 =
        d\frac{\left(1-\left(1+\frac{r}{k}\right)^{-kt}\right)}{\frac{r}{k}}$ \\
        \hline
        \end{tabular}
        \end{center}
        \end{minipage}}
}{% presentation
}{% nopresentation
}{% comments
}


\content{
        \begin{example} Suppose that I take out an auto loan from the bank. What
        does the formula for this situation look like? 
        \end{example}
}{% presentation
\vspace{3in}
}{% nopresentation
\vspace{3in}
}{% comments
}

\content{
        \begin{example} You want to take out a \$140,000 mortgage (home loan).
        The interest rate on the loan is 6\%,
        and the loan is for 30 years. How much will your monthly payments be?
        \end{example}
}{% presentation
\vspace{3in}
}{% nopresentation
\vspace{3in}
}{% comments
}

\content{
        \begin{example} Today's mortgage rates are around 6.38\%. Suppose you
        take out a \$250,000 mortgage. What will your monthly payment be to the
        bank?
        What is the total amount you will pay the bank over the next 30 years? 
        \end{example}
}{% presentation
\vspace{3in}
}{% nopresentation
\vspace{3in}
}{% comments
}

\content{
        \begin{example} Mortgage rates tend to be between 6.13\% and 7.26\%, and
        is affected by many things such as the borrower's credit score, the
        amount of the loan, inflation, etc. Suppose that you wish to take out a
        mortgage of \$300,000 from the bank of your choice. Compare the outcomes
        of the following two scenarios: 
        \begin{enumerate}
        \item You take out a 30 year loan, with payments made monthly, 
        during a year in which rates are slightly
        higher, at 6.85\%. What is the total amount you will pay the bank over
        the next 30 years? 
        \item You wait a year, and the rate you are offered now by the bank is
        6.43\%. The loan is still 30 years with payments made monthly. What is
        the total amount you will pay the bank over the next 30 years? 
        \end{enumerate}
        \end{example}
}{% presentation
\vspace{4in}
}{% nopresentation
\vspace{4in}
}{% comments
For the first, $d_1 = 1965.78$, and the second $d_2 = 1882.41$. Total amounts
paid are $A_1 = 707680.80$ and $A_2=677667.60$. That's a \$30013.20 difference! 
}

\content{
        \begin{example} Suppose that you have a 5 year auto loan for which you pay
        \$325 per month. If the interest rate is 5\%, how much will you still
        owe the bank after 3 years of making payments?  (It helps to think of
        this one from this perspective: If the bank invests money into you, and
        they wish to make \$325 `withdrawals' once a month for 2 years, how much
        do they need to have in the account). 
        \end{example}
}{% presentation
\vspace{3in}
}{% nopresentation
\vspace{3in}
}{% comments
}

\content{
        \begin{example} A couple purchases a home with a \$180,000 mortgage at
        4\% for 30 years with monthly payments. 
        What will the remaining balance on their mortgage be after 5 years?
        (hint: you first need to know what their monthly payment is, then it's
        just like the previous example)
        \end{example}
}{% presentation
\vspace{4in}
}{% nopresentation
\vspace{4in}
}{% comments
}
